\chapter{2008年辽宁卷（理）}

\section{单选题}

\begin{problem}
已知集合 \(M = \left\{x \mid \frac{x + 3}{x - 1} < 0\right\}\)，\(N = \{x \mid x \leqslant -3\}\)，则集合 \(\{x \mid x \geqslant 1\} =\)（\quad）

\choices
    {\(M \cap N\)}
    {\(M \cup N\)}
    {\(\complement_{\R}(M\cap N)\)}
    {\(\complement_{\R}(M \cup N)\)}
\end{problem}

\begin{answer}
D
\end{answer}

\begin{problem}
\(\displaystyle\lim_{n\to \infty}\frac{1 + 3 + 5 + \cdots + (2n - 1)}{n(2n + 1)} =\)（\quad）

\choices
    {\(\frac{1}{4}\)}
    {\(\frac{1}{2}\)}
    {1}
    {2}
\end{problem}

\begin{answer}
B
\end{answer}

\begin{problem}
圆 \(x^{2} + y^{2} = 1\) 与直线 \(y = kx + 2\) 没有公共点的充要条件是（\quad）

\choices
    {\(k \in (-\sqrt{2}, \sqrt{2})\)}
    {\(k \in (-\infty, -\sqrt{2}) \cup (\sqrt{2}, +\infty)\)}
    {\(k \in (-\sqrt{3}, \sqrt{3})\)}
    {\(k \in (-\infty, -\sqrt{3}] \cup (\sqrt{3}, +\infty)\)}
\end{problem}

\begin{answer}
C
\end{answer}

\begin{problem}
复数 \( \frac{1}{-2+\i} + \frac{1}{1-2\i} \) 的虚部是（\quad）

\choices
    {\(\frac{1}{5}\i\)}
    {\(\frac{1}{5}\)}
    {\(-\frac{1}{5}\i\)}
    {\(-\frac{1}{5}\)}
\end{problem}

\begin{answer}
B
\end{answer}

\begin{problem}
已知 \(O\)，\(A\)，\(B\) 是平面上的三个点，直线 \(AB\) 上有一点 \(C\)，满足 \(2\overrightarrow{AC} + \overrightarrow{CB} = 0\)，则 \(\overrightarrow{OC} =\)（\quad）

\choices
    {\(2\overrightarrow{OA} -\overrightarrow{OB}\)}
    {\(-\overrightarrow{OA} +2\overrightarrow{OB}\)}
    {\(\frac{2}{3}\overrightarrow{OA} -\frac{1}{3}\overrightarrow{OB}\)}
    {\(-\frac{1}{3}\overrightarrow{OA} +\frac{2}{3}\overrightarrow{OB}\)}
\end{problem}

\begin{answer}
A
\end{answer}

\begin{problem}
设 \(P\) 为曲线 \(C: y = x^2 + 2x + 3\) 上的点，且曲线 \(C\) 在点 \(P\) 处切线倾斜角的取值范围是 \(\left[0, \frac{\pi}{4}\right]\)，则点 \(P\) 横坐标的取值范围是（\quad）

\choices
    {\(\left[-1, -\frac{1}{2}\right]\)}
    {\([-1,0]\)}
    {[0,1]}
    {\(\left[\frac{1}{2}, 1\right]\)}
\end{problem}

\begin{answer}
A
\end{answer}

\begin{problem}
\(4\text{ 张}\)卡片上分别写有数字 \(1\)，\(2\)，\(3\)，\(4\)，从这 \(4\text{ 张}\)卡片中随机抽取 \(2\text{ 张}\)，则取出的 \(2\text{ 张}\)卡片上的数字之和为奇数的概率为（\quad）

\choices
    {\(\frac{1}{3}\)}
    {\( \frac{1}{2} \)}
    {\( \frac{2}{3} \)}
    {\(\frac{3}{4}\)}
\end{problem}

\begin{answer}
C
\end{answer}

\begin{problem}
将函数 \( y = 2^{x} + 1 \) 的图象按向量 \( \bs{a} \) 平移得到函数 \( y = 2^{x + 1} \) 的图象，则（\quad）

\choices
    {\(\bs{a} = (-1, - 1)\)}
    {\(\bs{a} = (1, - 1)\)}
    {\(\bs{a} = (1, 1)\)}
    {\(\bs{a} = (-1, 1)\)}
\end{problem}

\begin{answer}
A
\end{answer}

\begin{problem}
一生产过程有 \(4\text{ 道}\)工序，每道工序需要安排一人照看. 现从甲，乙，丙等 \(6\text{ 名}\)工人中安排 \(4\text{ 人}\)分别照看一道工序，第一道工序只能从甲，乙两工人中安排 \(1\text{ 人}\)，第四道工序只能从甲，丙两工人中安排 \(1\text{ 人}\)，则不同的安排方案共有（\quad）

\choices
    {\(24\text{ 种}\)}
    {\(36\text{ 种}\)}
    {\(48\text{ 种}\)}
    {\(72\text{ 种}\)}
\end{problem}

\begin{answer}
B
\end{answer}

\begin{problem}
已知点 \(P\) 是抛物线 \(y^{2} = 2x\) 上的一个动点，则点 \(P\) 到点 \((0, 2)\) 的距离与 \(P\) 到该抛物线准线的距离之和的最小值为（\quad）

\choices
    {\(\frac{\sqrt{17}}{2}\)}
    {3}
    {\(\sqrt{5}\)}
    {\(\frac{9}{2}\)}
\end{problem}

\begin{answer}
A
\end{answer}

\begin{problem}
在正方体 \(ABCD - A_{1}B_{1}C_{1}D_{1}\) 中，\(E\)，\(F\) 分别为棱 \(AA_{1}\)，\(CC_{1}\) 的中点，则在空间中与三条直线 \(A_{1}D_{1}\)，\(EF\)，\(CD\) 都相交的直线（\quad）

\choices
    {不存在}
    {有且只有两条}
    {有且只有三条}
    {有无数条}
\end{problem}

\begin{answer}
D
\end{answer}

\begin{problem}
设 \( f(x) \) 是连续的偶函数，且当 \( x > 0 \) 时 \( f(x) \) 是单调函数，则满足 \( f(x) = f\left(\frac{x + 3}{x + 4}\right) \) 的所有 \( x \) 之和为（\quad）

\choices
    {\(-3\)}
    {3}
    {\(-8\)}
    {8}
\end{problem}

\begin{answer}
C
\end{answer}

\section{填空题}

\begin{problem}
函数 \( y = \begin{cases}
x + 1, & x < 0,\\
\e^x, & x \geqslant 0,
\end{cases} \) 的反函数是\fillinblank{}.
\end{problem}

\begin{answer}
\(y = \begin{cases}
x - 1, & x < 1,\\
\ln x, & x \geqslant 1
\end{cases}\)
\end{answer}

\begin{problem}
在体积为 \(4\sqrt{3}\pi\) 的球的表面上有 \(A\)，\(B\)，\(C\) 三点，\(AB = 1\)，\(BC = \sqrt{2}\)，\(A\)，\(C\) 两点的球面距离为 \(\frac{\sqrt{3}}{3}\pi\)，则球心到平面 \(ABC\) 的距离为\fillinblank{}.
\end{problem}

\begin{answer}
\(\frac{3}{2}\)
\end{answer}

\begin{problem}
已知 \((1 + x + x^{2})\left(x + \frac{1}{x^{3}}\right)^{n}\) 的展开式中没有常数项，\(n \in \N^{*}\) 且 \(2 \leqslant n \leqslant 8\)，则 \(n =\)\fillinblank{}.
\end{problem}

\begin{answer}
\(5\)
\end{answer}

\begin{problem}
已知 \(f(x) = \sin \left( \omega x + \frac{\pi}{3} \right)\)（\(\omega > 0\)），\( f\left(\frac{\pi}{6}\right) = f\left(\frac{\pi}{3}\right) \)，且 \( f(x) \) 在区间 \( \left(\frac{\pi}{6}, \frac{\pi}{3}\right) \) 有最小值，无最大值，则 \( \omega = \)\fillinblank{}.
\end{problem}

\begin{answer}
\(\frac{14}{3}\)
\end{answer}

\section{解答题}

\begin{problem}
在 \(\triangle ABC\) 中，内角 \(A\)，\(B\)，\(C\) 对边的边长分别是 \(a\)，\(b\)，\(c\). 已知 \(c = 2\)，\(C = \frac{\pi}{3}\).

\begin{enumerate}
\item 若 \(\triangle ABC\) 的面积等于 \(\sqrt{3}\)，求 \(a\)，\(b\)；
\item 若 \(\sin C + \sin (B - A) = 2\sin 2A\)，求 \(\triangle ABC\) 的面积.
\end{enumerate}
\end{problem}

\begin{solution}
\begin{enumerate}
\item 由三角形面积公式，\(\frac{1}{2}ab\sin C=\sqrt{3}\)，即 \(\frac{\sqrt{3}}{4}ab=\sqrt{3}\)，所以 \(ab=4\). 由余弦定理，\(c^{2}=a^{2}+b^{2}-2ab\cos C=a^{2}+b^{2}-ab\)，故 \(a^{2}+b^{2}=8\). 从而 \((a+b)^{2}=a^{2}+b^{2}+2ab=16\)，由于 \(a,b>0\)，所以 \(a+b=4\). 因此 \(a,b\) 是方程 \(t^{2}-4t+4=0\) 的两个正根，得 \(a=b=2\).

\item 由 \(A+B+C=\pi\) 及正弦和差公式，\(\sin C+\sin(B-A)=2\sin\frac{B+C-A}{2}\cos\frac{C-B+A}{2}=2\cos A\sin B\). 所以题设条件等价于 \(2\cos A\sin B=2\sin 2A=4\sin A\cos A\).

若 \(\cos A=0\)，则 \(A=\frac{\pi}{2}\)，从而 \(B=\frac{\pi}{6}\). 由正弦定理，\(a=\frac{c\sin A}{\sin C}=\frac{4}{\sqrt{3}}\)，\(b=\frac{c\sin B}{\sin C}=\frac{2}{\sqrt{3}}\)，于是 \(S=\frac{1}{2}ab\sin C=\frac{2\sqrt{3}}{3}\).

若 \(\cos A\ne0\)，则 \(\sin B=2\sin A\). 又 \(B=\frac{2\pi}{3}-A\)，所以 \(\sin B=\sin\left(\frac{2\pi}{3}-A\right)=\frac{\sqrt{3}}{2}\cos A+\frac{1}{2}\sin A\). 因此 \(\sqrt{3}\cos A=3\sin A\)，即 \(\tan A=\frac{1}{\sqrt{3}}\). 因为 \(0<A<\frac{2\pi}{3}\)，得 \(A=\frac{\pi}{6}\)，从而 \(B=\frac{\pi}{2}\). 由正弦定理，\(a=\frac{2}{\sqrt{3}}\)，\(b=\frac{4}{\sqrt{3}}\)，所以此时同样有 \(S=\frac{1}{2}ab\sin C=\frac{2\sqrt{3}}{3}\). 故 \(\triangle ABC\) 的面积为 \(\frac{2\sqrt{3}}{3}\).
\end{enumerate}
\end{solution}

\begin{problem}
某批发市场对某种商品的周销售量（单位：吨）进行统计，最近 \(100\text{ 周}\)的统计结果如下表所示：

\begin{center}
\begin{tabular}{c|ccc}
周销售量 & \(2\) & \(3\) & \(4\) \\
\hline
频数 & \(20\) & \(50\) & \(30\)
\end{tabular}
\end{center}

\begin{enumerate}
\item 根据上面统计结果，求周销售量分别为 \(2\text{ 吨}\)，\(3\text{ 吨}\)和 \(4\text{ 吨}\)的频率；
\item 已知每吨该商品的销售利润为 \(2\text{ 千元}\)，\(\xi\) 表示该种商品两周销售利润的和（单位：千元）. 若以上述频率作为概率，且各周的销售量相互独立，求 \(\xi\) 的分布列和数学期望.
\end{enumerate}
\end{problem}

\begin{solution}
\begin{enumerate}
\item 由频率等于频数除以总周数，得周销售量为 \(2\) 吨、\(3\) 吨和 \(4\) 吨的频率分别为 \(\frac{20}{100}=0.2\)，\(\frac{50}{100}=0.5\)，\(\frac{30}{100}=0.3\).

\item 设一周的销售量为 \(X\)，则 \(P(X=2)=0.2\)，\(P(X=3)=0.5\)，\(P(X=4)=0.3\). 设两周销售量分别为 \(X_1\)，\(X_2\)，则 \(\xi=2X_1+2X_2\). 由各周销售量相互独立，\(\xi\) 的分布列如下：

\begin{center}
\begin{tabular}{c|ccccc}
\(\xi\) & \(8\) & \(10\) & \(12\) & \(14\) & \(16\) \\
\hline
\(P\) & \(0.04\) & \(0.20\) & \(0.37\) & \(0.30\) & \(0.09\)
\end{tabular}
\end{center}

其中 \(P(\xi=8)=0.2^{2}=0.04\)，\(P(\xi=10)=2\times0.2\times0.5=0.20\)，\(P(\xi=12)=2\times0.2\times0.3+0.5^{2}=0.37\)，\(P(\xi=14)=2\times0.5\times0.3=0.30\)，\(P(\xi=16)=0.3^{2}=0.09\). 因此 \(E\xi=8\times0.04+10\times0.20+12\times0.37+14\times0.30+16\times0.09=12.4=\frac{62}{5}\)，故 \(\xi\) 的数学期望为 \(\frac{62}{5}\) 千元.
\end{enumerate}
\end{solution}

\begin{problem}
如图，在棱长为 \(1\) 的正方体 \(ABCD - A'B'C'D'\) 中，\(AP = BQ = b\)（\(0 < b < 1\)），截面 \(PQEF \parallel A'D\)，截面 \(PQGH \parallel AD'\).

\begin{enumerate}
\item 证明：平面 \(PQEF\) 和平面 \(PQGH\) 互相垂直；
\item 证明：截面 \(PQEF\) 和截面 \(PQGH\) 面积之和是定值，并求出这个值；
\item 若 \(D'E\) 与平面 \(PQEF\) 所成的角为 \(45^{\circ}\)，求 \(D'E\) 与平面 \(PQGH\) 所成角的正弦值.
\end{enumerate}

\bitmapfigure[max width=0.34\linewidth]{img/2008/liaoning_science/q19_fig1.png}
\end{problem}

\begin{solution}
建立空间直角坐标系，取 \(A\) 为原点，令 \(A(0,0,0)\)，\(B(1,0,0)\)，\(D(0,1,0)\)，\(A'(0,0,1)\). 于是 \(D'(0,1,1)\)，且 \(P(0,0,b)\)，\(Q(1,0,b)\).
\begin{enumerate}
\item 有 \(\overrightarrow{PQ}=(1,0,0)\)，\(\overrightarrow{A'D}=(0,1,-1)\)，\(\overrightarrow{AD'}=(0,1,1)\). 因此，平面 \(PQEF\) 的一个法向量和平面 \(PQGH\) 的一个法向量可分别取为 \(\bs{n}_1=\overrightarrow{PQ}\times\overrightarrow{A'D}=(0,1,1)\)，\(\bs{n}_2=\overrightarrow{PQ}\times\overrightarrow{AD'}=(0,-1,1)\). 由于 \(\bs{n}_1\cdot\bs{n}_2=0\)，所以平面 \(PQEF\) 和平面 \(PQGH\) 互相垂直.

\item 从 \(P\) 沿方向 \(\overrightarrow{A'D}\) 移动，在线段 \(PQEF\) 内的参数范围为 \(0\leqslant t\leqslant b\)，所以该截面中与 \(PQ\) 垂直的边长为 \(b\left|A'D\right|=b\sqrt{2}\). 同理，从 \(P\) 沿方向 \(\overrightarrow{AD'}\) 移动，在线段 \(PQGH\) 内的参数范围为 \(0\leqslant t\leqslant1-b\)，所以该截面中与 \(PQ\) 垂直的边长为 \((1-b)\left|AD'\right|=(1-b)\sqrt{2}\). 因为 \(PQ\) 与上述两个方向均垂直，故 \(S_{PQEF}=1\cdot b\sqrt{2}=b\sqrt{2}\)，\(S_{PQGH}=1\cdot(1-b)\sqrt{2}=(1-b)\sqrt{2}\)，从而 \(S_{PQEF}+S_{PQGH}=\sqrt{2}\).

\item 平面 \(PQEF\) 内，点 \(E\) 的坐标为 \(E(1,b,0)\)，所以 \(\overrightarrow{D'E}=(1,b-1,-1)\). 记直线 \(D'E\) 与平面 \(PQEF\) 所成的角为 \(\theta\). 由直线与平面所成角的定义，\(\sin\theta=\frac{\left|\overrightarrow{D'E}\cdot\bs{n}_1\right|}{\left|\overrightarrow{D'E}\right|\left|\bs{n}_1\right|}=\frac{2-b}{\sqrt{2+(1-b)^2}\sqrt{2}}\). 由题意 \(\theta=45^\circ\)，故 \(\frac{2-b}{\sqrt{2+(1-b)^2}\sqrt{2}}=\frac{\sqrt{2}}{2}\). 两边平方并整理，得 \((2-b)^2=2+(1-b)^2\)，即 \(b=\frac{1}{2}\). 此时 \(\left|\overrightarrow{D'E}\right|=\sqrt{1+\frac{1}{4}+1}=\frac{3}{2}\). 记直线 \(D'E\) 与平面 \(PQGH\) 所成的角为 \(\varphi\)，则 \(\sin\varphi=\frac{\left|\overrightarrow{D'E}\cdot\bs{n}_2\right|}{\left|\overrightarrow{D'E}\right|\left|\bs{n}_2\right|}=\frac{\frac{1}{2}}{\frac{3}{2}\sqrt{2}}=\frac{\sqrt{2}}{6}\).
\end{enumerate}
\end{solution}

\begin{problem}
在直角坐标系 \(xOy\) 中，点 \(P\) 到两点 \((0, -\sqrt{3})\)，\((0, \sqrt{3})\) 的距离之和为 4，设点 \(P\) 的轨迹为 \(C\)，直线 \(y = kx + 1\) 与 \(C\) 交于 \(A\)，\(B\) 两点.

\begin{enumerate}
\item 写出 \(C\) 的方程；
\item 若 \(\overrightarrow{OA} \perp \overrightarrow{OB}\)，求 \(k\) 的值；
\item 若点 \(A\) 在第一象限，证明：当 \(k > 0\) 时，恒有 \(\left|\overrightarrow{OA}\right| > \left|\overrightarrow{OB}\right|\).
\end{enumerate}
\end{problem}

\begin{solution}
\begin{enumerate}
\item 两焦点为 \(F_1(0,-\sqrt{3})\)，\(F_2(0,\sqrt{3})\)，且到两焦点的距离之和为 \(4\)，所以 \(2a=4\)，即 \(a=2\). 焦距参数满足 \(c=\sqrt{3}\)，故 \(b^2=a^2-c^2=4-3=1\). 由于长轴在 \(y\) 轴上，轨迹 \(C\) 的方程为 \(\frac{x^2}{1}+\frac{y^2}{4}=1\)，即 \(4x^2+y^2=4\).

\item 设直线与椭圆的两个交点的横坐标分别为 \(x_1\)，\(x_2\)，则交点分别为 \(A=(x_1,kx_1+1)\)，\(B=(x_2,kx_2+1)\). 将 \(y=kx+1\) 代入 \(4x^2+y^2=4\)，得 \((k^2+4)x^2+2kx-3=0\). 所以由韦达定理，\(x_1+x_2=-\frac{2k}{k^2+4}\)，\(x_1x_2=-\frac{3}{k^2+4}\). 由 \(\overrightarrow{OA}\perp\overrightarrow{OB}\)，得 \(0=\overrightarrow{OA}\cdot\overrightarrow{OB}=x_1x_2+(kx_1+1)(kx_2+1)=(1+k^2)x_1x_2+k(x_1+x_2)+1\). 代入韦达定理的结果，得 \(0=1-\frac{3(1+k^2)+2k^2}{k^2+4}=\frac{1-4k^2}{k^2+4}\). 因此 \(k=\pm\frac{1}{2}\).

\item 设 \(x_A=u\)，\(x_B=v\). 因为点 \(A\) 在第一象限，所以 \(u>0\). 又 \(uv=-\frac{3}{k^2+4}<0\)，故 \(v<0\)，从而 \(u-v>0\). 由上面的韦达定理结果，\(u+v=-\frac{2k}{k^2+4}\). 于是
\[
\begin{gathered}
\left|\overrightarrow{OA}\right|^2-\left|\overrightarrow{OB}\right|^2=u^2+(ku+1)^2-v^2-(kv+1)^2\\
=(u-v)\bigl((1+k^2)(u+v)+2k\bigr)=(u-v)\left(-\frac{2k(1+k^2)}{k^2+4}+2k\right)\\
=\frac{6k(u-v)}{k^2+4}>0.
\end{gathered}
\]
所以 \(\left|\overrightarrow{OA}\right|>\left|\overrightarrow{OB}\right|\).
\end{enumerate}
\end{solution}

\begin{problem}
在数列 \(\{a_n\}\)，\(\{b_n\}\) 中，\(a_1 = 2\)，\(b_1 = 4\)，且 \(a_n\)，\(b_n\)，\(a_{n+1}\) 成等差数列，\(b_n\)，\(a_{n+1}\)，\(b_{n+1}\) 成等比数列.

\begin{enumerate}
\item 求 \(a_{2}\)，\(a_{3}\)，\(a_{4}\) 及 \(b_{2}\)，\(b_{3}\)，\(b_{4}\)，由此猜测 \(\{a_{n}\}\)，\(\{b_{n}\}\) 的通项公式，并证明你的结论；
\item 证明：\(\frac{1}{a_1 + b_1} + \frac{1}{a_2 + b_2} + \cdots + \frac{1}{a_n + b_n} < \frac{5}{12}\).
\end{enumerate}
\end{problem}

\begin{solution}
\begin{enumerate}
\item 由等差数列条件和等比数列条件，分别有 \(a_n+a_{n+1}=2b_n\)，\(a_{n+1}^{2}=b_nb_{n+1}\). 由 \(b_1=4\ne0\) 以及下面归纳得到的正性，递推时可写为 \(a_{n+1}=2b_n-a_n\)，\(b_{n+1}=\frac{a_{n+1}^{2}}{b_n}\). 先计算前几项：\(a_2=2b_1-a_1=6\)，\(b_2=\frac{a_2^2}{b_1}=9\)，\(a_3=2b_2-a_2=12\)，\(b_3=\frac{a_3^2}{b_2}=16\)，\(a_4=2b_3-a_3=20\)，\(b_4=\frac{a_4^2}{b_3}=25\). 由此猜测 \(a_n=n(n+1)\)，\(b_n=(n+1)^2\). 下面用数学归纳法证明：当 \(n=1\) 时，\(a_1=2=1\times2\)，\(b_1=4=2^2\)，结论成立. 假设 \(a_n=n(n+1)\)，\(b_n=(n+1)^2\)，则 \(b_n>0\)，并且 \(a_{n+1}=2(n+1)^2-n(n+1)=(n+1)(n+2)>0\)，\(b_{n+1}=\frac{a_{n+1}^{2}}{b_n}=\frac{[(n+1)(n+2)]^2}{(n+1)^2}=(n+2)^2\). 所以对任意正整数 \(n\)，都有 \(a_n=n(n+1)\)，\(b_n=(n+1)^2\).

\item 由通项公式，\(a_i+b_i=i(i+1)+(i+1)^2=(i+1)(2i+1)\). 当 \(n=1\) 时，左边等于 \(\frac{1}{6}<\frac{5}{12}\). 当 \(n\geqslant2\) 时，对 \(i\geqslant2\)，有 \(2i+1>2i\)，从而
\[
\begin{gathered}
\sum_{i=1}^{n}\frac{1}{a_i+b_i}=\frac{1}{6}+\sum_{i=2}^{n}\frac{1}{(i+1)(2i+1)}<\frac{1}{6}+\frac{1}{2}\sum_{i=2}^{n}\frac{1}{i(i+1)}\\
=\frac{1}{6}+\frac{1}{2}\sum_{i=2}^{n}\left(\frac{1}{i}-\frac{1}{i+1}\right)=\frac{1}{6}+\frac{1}{2}\left(\frac{1}{2}-\frac{1}{n+1}\right)<\frac{5}{12}.
\end{gathered}
\]
故不等式得证.
\end{enumerate}
\end{solution}

\begin{problem}
设函数 \( f(x) = \frac{\ln x}{1 + x} - \ln x + \ln (x + 1) \).

\begin{enumerate}
\item 求 \( f(x) \) 的单调区间和极值.
\item 是否存在实数 \(a\)，使得关于 \(x\) 的不等式 \(f(x) \geqslant a\) 的解集为 \((0, +\infty)\)？若存在，求 \(a\) 的取值范围；若不存在，试说明理由.
\end{enumerate}
\end{problem}

\begin{solution}
\begin{enumerate}
\item 函数定义域为 \(x>0\)，且 \(f(x)=\ln(x+1)-\frac{x}{x+1}\ln x\). 求导得 \(f'(x)=\frac{1}{x+1}-\left(\frac{\ln x}{(x+1)^2}+\frac{1}{x+1}\right)=-\frac{\ln x}{(x+1)^2}\). 因此，当 \(0<x<1\) 时，\(f'(x)>0\)；当 \(x>1\) 时，\(f'(x)<0\). 所以 \(f(x)\) 在 \((0,1)\) 上单调递增，在 \((1,+\infty)\) 上单调递减，并且 \(f(1)=\ln2\). 又因为 \(\lim_{x\to0^+}x\ln x=0\)，所以 \(\lim_{x\to0^+}f(x)=0\). 将函数改写为 \(f(x)=\ln\left(1+\frac{1}{x}\right)+\frac{\ln x}{x+1}\)，可得 \(\lim_{x\to+\infty}f(x)=0\). 所以 \(f(x)\) 的最大值为 \(\ln2\)，在 \(x=1\) 时取得；函数没有最小值，但下确界为 \(0\).

\item 由单调性和两个端点极限，\(f(x)>0\) 对一切 \(x>0\) 成立. 因此当 \(a\leqslant0\) 时，不等式 \(f(x)\geqslant a\) 的解集为 \((0,+\infty)\). 若 \(a>0\)，由 \(\lim_{x\to0^+}f(x)=0\)，可取充分接近 \(0\) 的正数 \(x\)，使 \(f(x)<a\)，此时解集不可能是整个 \((0,+\infty)\). 故存在这样的实数 \(a\)，其取值范围为 \(a\leqslant0\).
\end{enumerate}
\end{solution}
